\documentclass[twoside,reqno]{amsart}
    \usepackage[all]{xy}
    \usepackage{graphicx,amsmath,amssymb,amsthm,amsfonts,amscd,euscript,tikz}
    \usepackage{fancyhdr,hyperref,booktabs,xspace,mathrsfs,pdflscape}
    \usepackage[a4paper,left=3cm,right=3cm,top=3.5cm,bottom=3.5cm]{geometry}
    \usepackage{float}
    \usepackage{subcaption}
    \usepackage{xepersian}
    \settextfont[Scale=1.5]{IRNazanin}
    \setdigitfont[Scale=1.3]{IRNazanin}
    \setlatintextfont[Scale=1.2]{Times New Roman}
    \emergencystretch=3em
    \linespread{1.8}
    \renewcommand{\refname}{مراجع}
    \csname@twosidetrue\endcsname
    \pagestyle{fancy}
    \fancyhf{} 
    \fancyhead[OL,EL]{}
    \fancyhead[OR]{\small\rightmark}
    \fancyhead[ER]{\small\leftmark}
    \fancyfoot[C]{\thepage}
    \newcommand{\Lfootnote}[1]{\hspace{-1mm}\LTRfootnote{\rm #1}~}

    \begin{document}
    \thispagestyle{empty}

    \pagestyle{myheadings}
    \def\n{\noindent}%

    \numberwithin{equation}{section}
    \def\mapdown#1{\Big\downarrow\rlap
    {$\vcenter{\hbox{$\scriptstyle#1$}}$}}
    \bibliographystyle{amsplain}

    \setcounter{page}{1} \noindent


    %%%%%%%%%%%%%%%%%%%%%%%%%%%%%%%%%%%%%%%
    \vspace{1.5cm}
    \begin{center}
    {\Large \bf تحلیل داده‌های مسکن}\\
    \vspace{0.5cm}

    \textsuperscript{1,2}\lr{Hoda Bornak}\\
    \vspace{0.1cm}
    \textsuperscript{1} \lr{Semnan University}\\
    \vspace{0.1cm}
    \textsuperscript{2} \texttt{hodabornak@semnan.ac.ir}

    \end{center}
    %%%%%%%%%%%%%%%%%
    \vspace*{0.5cm}
    \begin{quotation}
    \noindent
    {\footnotesize
    \textbf{چکیده.}
 داده‌های مورد استفاده در این پژوهش شامل اطلاعات مربوط به معاملات مسکن در سطح کشور است. این داده‌ها شامل متغیرهایی نظیر قیمت هر متر مربع \lr{gheimat}، متراژ \lr{masahat}، سن بنا \lr{omre bana}، استان \lr{State}، شهر \lr{City}، نوع سازه \lr{skeleton} و  منطقه \lr{mantaghe} می‌باشند.
پس از حذف رکوردهای تکراری و اعمال فیلترها، تعداد کل رکوردهای معتبر به ۱۸۲۹۲ رسید که شامل ۳۰ استان و  ۲۴۸ شهر مختلف می‌باشد. همچنین داده‌های مربوط به شهر تهران که شامل ۶۱۱۵ رکورد از ۲۲ منطقه شهری است، به‌عنوان زیرمجموعه‌ای جداسازی شد تا تحلیل‌های دقیق‌تری بر روی پایتخت انجام شود.
ابتدا داده‌ها پاکسازی و توصیف شده و سپس با استفاده از تکنیک‌های بصری‌سازی  و مدل جنگل تصادفی، الگوهای حاکم بر بازار مسکن بررسی شده است. نتایج نشان می‌دهد که عامل «شهر» با ضریب اهمیت ۵/۵۰ درصد، مهم‌ترین عامل تعیین‌کننده قیمت مسکن است و پس از آن به ترتیب «متراژ» (۹/۲۷ درصد)، «روز معامله» (۶/۹ درصد)، «سن بنا» (۹/۸ درصد) و «نوع سازه» (۳/۲ درصد) قرار دارند. مدل پیش‌بینی با دقت ۹۴ درصد (\lr{$R^2 = 0.9429$}) توانسته است تغییرات قیمت را تبیین کند.
 
    }
    \end{quotation}
    
\section{مقدمه}
بازار مسکن به عنوان یکی از بخش‌های حیاتی اقتصاد هر کشور، همواره تحت تأثیر عوامل متعددی قرار دارد. درک عمیق این عوامل و نحوه تأثیرگذاری آن‌ها بر قیمت مسکن، برای سیاست‌گذاران، سرمایه‌گذاران و شهروندان از اهمیت بالایی برخوردار است. هدف این مطالعه، بررسی داده‌های واقعی بازار مسکن ایران به منظور شناسایی عوامل کلیدی مؤثر بر قیمت، و تحلیل تفاوت‌های منطقه‌ای و پیش‌بینی قیمت با استفاده از روش‌های یادگیری ماشین است.
\section{داده‌ها و روش‌شناسی}
\subsection{منبع و ساختار داده}
داده‌های استفاده شده در این مطالعه، از فایل \lr{Iran Maskan.xlsx} استخراج شده است. مجموعه داده پس از پاکسازی، شامل ۱۸۲۹۲ رکورد معتبر می‌باشد که ویژگی‌های زیر را دارد:
\begin{itemize}
    \item \textbf{متغیر هدف}: قیمت هر متر مربع مسکن ($gheimat$)
    \item \textbf{متغیرهای مستقل}: 
    \begin{itemize}
        \item متراژ واحد مسکونی ($masahat$)
        \item سن بنا ($omre bana$)
        \item نوع سازه ($skeleton$)
        \item استان ($State$)
        \item شهر ($City$)
        \item روز معامله ($day$)
        \item منطقه شهرداری در تهران ($mantaghe$)

    \end{itemize}
\end{itemize}

\subsection{پیش‌پردازش داده‌ها}
 مراحل زیر برای آماده‌سازی داده‌ها برای مدل‌سازی انجام شد:
\begin{enumerate}
    \item \textbf{پالایش اولیه}: رکوردهای تکراری و رکوردهای با قیمت صفر یا منفی،  متراژ کمتر از ۵ متر مربع و سن بنای منفی حذف شدند.
    \item \textbf{کدگذاری متغیرهای کیفی}: متغیرهای دسته‌ای (استان، شهر، منطقه و نوع سازه) با استفاده از تکنیک \lr{Label Encoding} به مقادیر عددی تبدیل شدند.
\end{enumerate}

\begin{table}
\caption{آمار توصیفی متغیرهای عددی}
\begin{tabular}{|c|c|c|c|}
\hline
& متراژ & قیمت هر متر مربع & سن بنا \\
\hline
میانگین & ۱۸۴ &  ۹۶۲,۶۶  & ۹ \\
انحراف‌معیار &  ۷۲۶,۲   & ۵۴۲,۶۹  & ۸ \\
حداقل & ۷ & ۱ & ۰ \\
چارک اول & ۶۵ & ۰۰۰,۲۰ & ۳ \\
میانه & ۸۲ & ۲۵۰,۴۱ & ۸ \\
چارک سوم & ۱۰۷ &  ۴۱۳,۸۸  & ۱۳ \\
حداکثر & ۰۳۹,۲۰۴ & ۴۲۰,۶۶۲  & ۹۸ \\
\hline
\end{tabular}
\label{tab1}
\end{table}

 جدول \ref{tab1} خلاصه‌ای از آمار داده‌ها را ارائه می‌دهد. همان‌طور که مشاهده می‌شود، میانگین قیمت هر متر مربع  حدود ۶۷۰۰۰ تومان است. انحراف‌معیار بالا (حدود  ۷۰۰۰۰ تومان) نشان‌دهنده پراکندگی قابل‌توجه قیمت‌ها در سطح کشور است. این موضوع حاکی از وجود تفاوت‌های قیمتی چشمگیر بین استان‌ها و شهرهای مختلف می‌باشد. میانگین سن بنا حدود ۹ سال است و بیشترین سن بنا به ۹۸ سال می‌رسد که نشان‌دهنده وجود بناهای قدیمی در مجموعه داده است.

\section{تحلیل اکتشافی و بصری‌سازی داده‌ها}
\subsection{تحلیل توزیع و پراکندگی قیمت}

در این بخش، توزیع قیمت هر متر مربع و پراکندگی آن در گروه‌های مختلف سن بنا و نوع سازه بررسی شده است.

\begin{figure}[H]
\centering
\includegraphics[width=0.75\textwidth]{1.png}
\vspace{-3mm}
\caption{هیستوگرام توزیع قیمت هر متر مربع}
\label{1}
\end{figure}

هیستوگرام شکل \ref{1} توزیع قیمت هر متر مربع را پس از پالایش اولیه داده‌ها نمایش می‌دهد. توزیع داده‌ها دارای چولگی مثبت است؛ به این معنی که اکثر واحدهای مسکونی در بازه‌های قیمتی پایین تا متوسط قرار دارند و تعداد محدودی از املاک با قیمت‌های بسیار بالا، دنباله‌ای کشیده در سمت راست توزیع ایجاد کرده‌اند. میانگین قیمت حدود ۶۷۰۰۰ تومان و میانه قیمت حدود ۴۱۰۰۰ تومان است. بزرگ‌تر بودن میانگین نسبت به میانه نشان می‌دهد که این مقادیر بسیار بزرگ و دورافتاده، موجب افزایش مقدار میانگین و عدم تقارن توزیع شده‌اند.

\begin{figure}[H]
\centering
\includegraphics[width=0.7\textwidth]{2.png}
\vspace{-4mm}
\caption{نمودار جعبه‌ای قیمت هر متر مربع بر اساس نوع سازه}
\label{2}
\end{figure}

نمودار جعبه‌ای شکل \ref{2} توزیع قیمت هر متر مربع را بر اساس نوع سازه نمایش می‌دهد. بر اساس این نمودار، الگوی توزیع قیمت در انواع مختلف سازه یکسان نیست. سازه‌های «بتن ـ فلز» و «بتنی» دارای میانه قیمت بالاتری نسبت به سازه‌های «فلزی» و «سایر» هستند، در حالی که اختلاف بین میانه این دو نوع سازه ناچیز است. گروه «سایر» که شامل سازه‌های «آجر بلوک سیمانی»، «خشتی ـ گلی»، «چوبی» و «بدون اسکلت» است، پایین‌ترین میانه و کمترین پراکندگی قیمت را دارد. در تمامی انواع سازه، تعداد زیادی مقادیر دورافتاده در بخش بالایی نمودار مشاهده می‌شود که نشان‌دهنده وجود واحدهایی با قیمت بسیار بالاتر از بازه متداول قیمت در بازار مسکن است. این موضوع نشان می‌دهد که حتی در میان املاک با یک نوع سازه مشخص نیز تنوع زیادی در قیمت هر متر مربع وجود دارد.

\vspace{5mm}

\begin{figure}[H]
\centering
\includegraphics[width=0.7\textwidth]{3.png}
\vspace{-3mm}
\caption{نمودار جعبه‌ای قیمت هر متر مربع بر اساس گروه سنی بنا}
\label{3}
\end{figure}

نمودار جعبه‌ای شکل \ref{3} توزیع قیمت هر متر مربع را در گروه‌های مختلف سنی نمایش می‌دهد. بر اساس این نمودار، میانه و پراکندگی قیمت در گروه‌های مختلف سنی متفاوت است. گروه «بیش از ۳۰ سال» بالاترین میانه و بیشترین پراکندگی قیمت را دارد. در تمامی گروه‌های سنی، مقادیر دورافتاده در بخش بالایی نمودار مشاهده می‌شود که نشان‌دهنده وجود واحدهایی با قیمت بسیار بالاتر از بازه متداول قیمت در بازار مسکن می‌باشد. دامنه این مقادیر دورافتاده در گروه‌های سنی «۰ تا ۵ سال» و  «۲۱ تا ۳۰ سال» بیشتر از سایر گروه‌ها است و در برخی از معاملات به مقادیر بیشتر از  ۰۰۰,۶۰۰ تومان نیز می‌رسد. 

\vspace{5mm}

\subsection{توزیع داده‌ها بر اساس استان}

در این بخش، ابتدا سهم هر استان از کل داده‌ها و سپس تعداد معاملات ثبت‌شده در هر استان بررسی شده است.

\begin{figure}[H]
\centering
\includegraphics[width=0.69\textwidth]{4.png}
\vspace{-10mm}
\caption{توزیع داده‌های املاک بر اساس استان}
\label{4}
\end{figure}

\begin{figure}[H]
\centering
\includegraphics[width=0.89\textwidth]{5.png}
\vspace{-3mm}
\caption{نمودار میله‌ای تعداد معاملات ثبت‌شده در همه استان‌ها}
\label{5}
\end{figure}

نمودار دایره‌ای شکل \ref{4} توزیع رکوردهای داده را بر اساس استان‌های مختلف نشان می‌دهد. استان تهران بیش از یک‌سوم کل داده‌ها را به خود اختصاص داده است. استان‌های البرز، فارس، خوزستان و گیلان در رتبه‌های بعدی قرار دارند. این تمرکز بالای داده‌ها در استان تهران نشان‌دهنده حجم بالای معاملات مسکن در پایتخت و همچنین دسترسی بیشتر به داده‌های این استان است.

نمودار میله‌ای شکل \ref{5} تعداد معاملات ثبت‌شده را در تمامی استان‌های کشور نمایش می‌دهد. پس از دو استان نخست، تعداد معاملات در سایر استان‌ها با شیب قابل‌توجهی کاهش می‌یابد. سهم استان‌هایی مانند کهگیلویه و بویراحمد، یزد، کرمان و سیستان و بلوچستان بسیار ناچیز است. این موضوع بیانگر آن است که اختلاف محسوسی در حجم معاملات استان‌های مختلف وجود دارد.


\subsection{میانگین قیمت هر متر مربع در استان‌های کشور}

در این بخش، میانگین قیمت هر متر مربع به‌صورت نمودار دوناتی برای استان‌های گران‌قیمت و نمودار میله‌ای برای تمامی استان‌ها نمایش داده شده است.

\begin{figure}[H]
\centering
\includegraphics[width=0.65\textwidth]{6.png}
\vspace{-10mm}
\caption{نمودار دوناتی میانگین قیمت هر متر مربع در استان‌های گران قیمت}
\label{6}
\end{figure} 

\begin{figure}[H]
\centering
\includegraphics[width=0.88\textwidth]{7.png}
\vspace{-3mm}
\caption{نمودار میله‌ای میانگین قیمت هر متر مربع در همه استان‌ها}
\label{7}
\end{figure}

نمودار دوناتی شکل \ref{6} نشان می‌دهد که بازار مسکن ایران با نابرابری‌های جغرافیایی شدیدی 
مواجه است. تهران به‌عنوان مرکز سیاسی و اقتصادی کشور، قیمت‌های بسیار بالاتری نسبت به سایر استان‌ها دارد. این اختلاف می‌تواند پیامدهای اجتماعی و اقتصادی مهمی مانند مهاجرت معکوس،  افزایش هزینه‌های زندگی و کاهش قدرت خرید خانوارها را به همراه داشته باشد.

\vspace{0.5mm}

نمودار  میله‌ای شکل \ref{7}  میانگین قیمت هر متر مربع را در تمامی استان‌ها نشان می‌دهد. بر اساس این نمودار، میانگین قیمت در تهران بیش از دو برابر میانگین قیمت در استان دوم (فارس) است. استان‌های خراسان رضوی، قم و اصفهان در رده‌های بعدی قرار دارند. در انتهای نمودار، استان‌های کهگیلویه و بویراحمد، لرستان و یزد پایین‌ترین میانگین قیمت را دارند. این شکاف عمیق نشان‌دهنده تمرکز شدید فعالیت‌های اقتصادی و تقاضای مسکن در تهران است.

\subsection{تحلیل الگوهای زمانی روزانه}

در این بخش، روند تعداد معاملات و  میانگین قیمت هر متر مربع  در روزهای خردادماه ۱۳۹۸ ترسیم شده است.

\begin{figure}[H]
\centering
\includegraphics[width=0.88\textwidth]{8.png}
\vspace{-3mm}
\caption{روند روزانه تعداد معاملات}
\label{8}
\end{figure}

\begin{figure}[H]
\centering
\includegraphics[width=0.88\textwidth]{9.png}
\vspace{-3mm}
\caption{روند روزانه میانگین قیمت هر متر مربع}
\label{9}
\end{figure}

نمودار خطی شکل \ref{8} روند تعداد معاملات ثبت‌شده را نشان می‌دهد. حجم معاملات در طول ماه دارای نوسانات قابل‌توجهی است. بیشترین تعداد معاملات در برخی از روزهای ابتدایی و اواسط ماه ثبت شده‌اند، در حالی که روزهای ۱۴ تا ۱۷، ۲۴ و ۳۱ خرداد کمترین میزان معاملات را به خود اختصاص داده‌اند.

\vspace{0.5mm}

نمودار خطی شکل \ref{9} روند روزانه میانگین قیمت هر متر مربع را نشان می‌دهد. بر اساس این نمودار، بیشترین و کمترین میانگین قیمت هر متر مربع در طول ماه به ترتیب مربوط به روزهای هفتم و چهاردهم خرداد هستند. این نوسانات روزانه می‌تواند نشان‌دهنده تأثیر عواملی مانند روزهای تعطیل، زمان وصول حقوق و سایر عوامل مؤثر بر بازار مسکن ‌باشد.

\subsection{بررسی قیمت مسکن در مناطق شهر تهران}

با توجه به اهمیت و سهم بالای داده‌های شهر تهران، تحلیل دقیق‌تری بر روی مناطق ۲۲ گانه این شهر انجام شده است تا تفاوت قیمت میان مناطق مختلف مشخص شود. 

\begin{figure}[H]
\centering
\includegraphics[width=0.88\textwidth]{10.png}
\vspace{-3mm}
\caption{نمودار میله‌ای قیمت هر متر مربع در همه مناطق شهر تهران}
\label{10}
\end{figure}

نمودار میله‌ای شکل \ref{10} میانگین قیمت هر متر مربع را در مناطق مختلف نمایش می‌دهد. نتایج نشان می‌دهد که مناطق شمالی و مرکزی تهران (مناطق ۱، ۲، ۳ و ۶) با اختلاف قابل‌توجهی گران‌ترین مناطق شهر هستند. منطقه ۳ با میانگین قیمت ۲۳۰۱۸۵ تومان گران‌ترین و منطقه ۲۰ با ۵۸۰۱۹ تومان ارزان‌ترین منطقه تهران است. روند کلی قیمت‌ها نشان می‌دهد که با حرکت از شمال به سمت جنوب و حاشیه شهر، قیمت مسکن به‌طور نسبی کاهش می‌یابد. این الگو با ساختار اقتصادی ـ اجتماعی شهر تهران و توزیع امکانات و خدمات شهری همخوانی دارد. اختلاف قیمت بین گران‌ترین و ارزان‌ترین منطقه به حدود ۱۷۲۰۰۰ تومان می‌رسد که نشان‌دهنده نابرابری شدید قیمت مسکن در پایتخت است.

\vspace{5mm}

\subsection{بررسی تأثیر سن بنا و نوع سازه بر قیمت در شهر تهران}

در این بخش، میانگین قیمت هر متر مربع بر اساس دو ویژگی سن بنا و نوع سازه در شهر تهران مورد بررسی قرار گرفته است.

\begin{figure}[H]
\centering
\includegraphics[width=0.8\textwidth]{11.png}
\vspace{-3mm}
\caption{نمودار میله‌ای میانگین قیمت بر اساس سن بنا در شهر تهران}
\label{11}
\end{figure}

\begin{figure}[H]
\centering
\includegraphics[width=0.8\textwidth]{12.png}
\vspace{-3mm}
\caption{نمودار میله‌ای میانگین قیمت بر اساس نوع سازه در شهر تهران}
\label{12}
\end{figure}

نمودار میله‌ای شکل \ref{11} میانگین قیمت را بر اساس سن بنا نشان می‌دهد. واحدهای نوساز گران‌ترین واحدها هستند. با افزایش سن بنا، میانگین قیمت تا رده سنی ۱۱ تا ۲۰ سال کاهش یافته و پس از آن، در گروه‌های سنی بالاتر، دوباره روندی افزایشی پیدا می‌کند. این موضوع می‌تواند ناشی از قرارگیری بخشی از این واحدها در مناطق قدیمی اما ارزشمند شهر تهران باشد.

\vspace{0.5mm}

نمودار میله‌ای شکل \ref{12} میانگین قیمت را بر اساس نوع سازه نمایش می‌دهد. سازه‌های «بتن ـ فلز» و «بتنی» به‌ترتیب با میانگین ۱۴۳۴۸۶ و ۱۴۰۸۳۵ تومان بالاترین قیمت را دارند و سازه‌های «فلزی» و «سایر» ارزان‌تر هستند. این تفاوت می‌تواند ناشی از مقاومت بیشتر، عمر مفید بالاتر و کیفیت ساخت بهتر سازه‌های بتنی و ترکیبی باشد.

\subsection{بررسی رابطه متراژ و سن بنا با قیمت کل}

در این بخش، به‌منظور بررسی رابطه قیمت کل ملک با متغیرهای متراژ و سن بنا، دو نمودار پراکندگی ترسیم شده‌اند.

\begin{figure}[H]
\centering
\includegraphics[width=0.65\textwidth]{13.png}
\vspace{-3mm}
\caption{نمودار پراکندگی رابطه متراژ و قیمت کل}
\label{13}
\end{figure}

\begin{figure}[H]
\centering
\includegraphics[width=0.65\textwidth]{14.png}
\vspace{-3mm}
\caption{نمودار پراکندگی رابطه سن بنا و قیمت کل}
\label{14}
\end{figure}

نمودار پراکندگی شکل \ref{13} رابطه بین متراژ و قیمت کل را نشان می‌دهد. به‌طور کلی، با افزایش متراژ، قیمت کل نیز افزایش می‌یابد، اما پراکندگی قابل‌توجه نقاط به‌ویژه در متراژهای بالاتر، نشان می‌دهد که متراژ به‌تنهایی نمی‌تواند تغییرات قیمت را به‌طور کامل توضیح دهد و عوامل دیگری نیز در تعیین قیمت نقش دارند.

\vspace{0.5mm}

نمودار پراکندگی شکل \ref{14}  رابطه بین سن بنا و قیمت کل را نمایش می‌دهد. برخلاف متراژ، الگوی مشخصی بین این دو متغیر مشاهده نمی‌شود و قیمت در تمامی گروه‌های سنی دارای پراکندگی بالایی است. این موضوع نشان می‌دهد که سایر ویژگی‌های ملک (مانند موقعیت مکانی)، نقش پررنگ‌تری در تعیین قیمت دارند.

\subsection{ماتریس همبستگی متغیرهای تأثیرگذار بر قیمت املاک}

در این بخش، شدت و جهت رابطه خطی بین متغیرها با استفاده از ماتریس همبستگی بررسی شده است.

\begin{figure}[H]
\centering
\includegraphics[width=0.6\textwidth]{15.png}
\caption{ماتریس همبستگی متغیرهای تأثیرگذار بر قیمت املاک}
\label{15}
\end{figure}

ماتریس همبستگی شکل \ref{15}، میزان همبستگی بین متغیرها را نشان می‌دهد. بیشترین همبستگی بین «قیمت کل» و «قیمت هر متر مربع» مشاهده می‌شود (با ضریب ۸۱۷/۰). این منطقی است؛ زیرا هرچه قیمت هر متر مربع بالاتر باشد، قیمت کل ملک نیز به طور قابل توجهی افزایش می‌یابد. همبستگی متراژ با قیمت کل (با ضریب ۰۲۷/۰) و متراژ با قیمت هر متر مربع (با ضریب ۰۳۰/۰-) بسیار ضعیف است. این موضوع نشان می‌دهد که در این داده‌ها، بزرگ‌تر بودن متراژ لزوماً به معنای بالاتر بودن قیمت کل یا قیمت هر متر مربع نیست (شاید به دلیل تنوع زیاد مناطق یا کیفیت متفاوت ساختمان‌ها باشد). سن بنا با سایر متغیرها همبستگی بسیار ناچیزی دارد. برای مثال، همبستگی آن با قیمت کل ۰۴۷/۰- است که نشان می‌دهد با افزایش سن بنا، قیمت کل کمی کاهش می‌یابد، اما این تأثیر در این مجموعه داده بسیار ضعیف است. متغیر روز ماه نیز تقریباً با هیچ‌کدام از متغیرهای دیگر همبستگی معناداری ندارد (اعداد نزدیک به صفر). این یعنی زمان انجام معامله در طول یک ماه، تأثیری بر قیمت یا متراژ ملک‌های معامله شده نداشته است.

\subsection{نمودار تعاملی تغییر قیمت در استان‌ها}

به منظور بررسی دقیق‌تر و پویای روند تغییر قیمت هر متر مربع در استان‌های مختلف طی روزهای ماه، نمودار تعاملی شکل \ref{16} با استفاده از کتابخانه Plotly در پایتون پیاده‌سازی شده است. این نمودار به کاربر امکان می‌دهد تا با استفاده از منوی کشویی، استان موردنظر خود را انتخاب کرده و روند تغییر میانگین قیمت را در آن استان مشاهده کند. 

\begin{figure}[H]
\centering
\includegraphics[width=0.9\textwidth]{16.png}
\caption{نمودار تعاملی تغییر قیمت در ‌استان‌ها}
\label{16}
\end{figure}

مزایای اصلی این نمودار تعاملی عبارتند از:

\begin{itemize}
    \item \textbf{نمایش اطلاعات دقیق:} با قرارگیری نشانگر ماوس روی هر نقطه، اطلاعات مربوط به روز معامله و میانگین قیمت هر متر مربع به‌صورت پاپ‌آپ نمایش داده می‌شود.
            \item \textbf{قابلیت بزرگ‌نمایی:} کاربر می‌تواند بخش‌های موردنظر خود را بزرگ‌نمایی کرده و جزئیات بیشتری را مشاهده کند.
               \item \textbf{قابلیت جابه‌جایی:} کاربر می‌تواند نمودار را جابه‌جا کند و در صورت نیاز با یک کلیک بر آیکون «خانه»، آن را به حالت اولیه بازگرداند.
                    \item \textbf{قابلیت ذخیره‌سازی:} نمودار را می‌توان به‌صورت فایل HTML ذخیره کرد و بدون نیاز به محیط پایتون، به‌صورت مستقل در مرورگر مشاهده کرد و به اشتراک گذاشت.
\end{itemize}

این نمودار تعاملی به‌ویژه برای ارائه‌های مدیریتی و جلسات تصمیم‌گیری که نیاز به بررسی دقیق و پویای داده‌ها دارند، بسیار مفید است.

\section{مدل‌سازی و تحلیل اهمیت ویژگی‌ها}
\subsection{انتخاب مدل جنگل تصادفی}

برای پیش‌بینی قیمت و تحلیل اهمیت ویژگی‌ها، از الگوریتم Regressor Forest Random استفاده شده است. انتخاب این الگوریتم به دلیل توانایی بالای آن در کار با داده‌های با ابعاد بالا، عدم حساسیت به مقیاس متغیرها، قابلیت مدل‌سازی روابط غیرخطی بین متغیرها و امکان ارائه میزان اهمیت هر ویژگی می‌باشد. 

\subsection{آموزش و ارزیابی مدل}

مدل Forest Random با ۱۰۰ درخت بر روی داده‌ها آموزش داده شد. معیار ارزیابی عملکرد مدل، ضریب تعیین ($R^2$) است که نشان می‌دهد مدل تا چه حد توانسته تغییرات متغیر هدف را توضیح دهد. مقدار $R^2$ محاسبه‌شده  برابر با ۹۴۲۹/۰ شده است که نشان‌دهنده دقت بالای مدل در پیش‌بینی قیمت مسکن بر اساس ویژگی‌های موجود می‌باشد.   

\subsection{تحلیل اهمیت ویژگی‌ها}

پس از آموزش مدل، میزان اهمیت هر یک از ویژگی‌های ورودی در تعیین قیمت مسکن استخراج شد. نتایج این تحلیل، در جدول \ref{tab:importance} و شکل \ref{17} نمایش داده شده است. هدف از این تحلیل، شناسایی میزان تأثیر هر یک از متغیرها بر قیمت نهایی است. 

\begin{figure}[H]
    \centering
    \includegraphics[width=0.7\textwidth]{17.png}
    \caption{نمودار اهمیت ویژگی‌های مختلف در تعیین قیمت مسکن}
    \label{17}
\end{figure}

\begin{table}[H]
\centering
\caption{اهمیت ویژگی‌های مدل در پیش‌بینی قیمت مسکن}
\label{tab:importance}
\begin{tabular}{lc}
\toprule
ویژگی & ضریب اهمیت \\
\midrule
شهر & ۵/۵۰\% \\
متراژ & ۹/۲۷\% \\
روز معامله & ۶/۹\% \\
سن بنا & ۹/۸\% \\
نوع سازه & ۳/۲\% \\
استان & ۸/۰\% \\
\bottomrule
\end{tabular}
\end{table}

این نتایج نشان می‌دهد که متغیر «شهر»، مهم‌ترین عامل تعیین‌کننده قیمت مسکن در ایران است. پس از آن، ویژگی‌های متراژ، روز معامله ملک و سن بنا در اولویت‌های بعدی قرار دارند.

\section{نتیجه‌گیری و جمع‌بندی}

در این پژوهش، داده‌های معاملات مسکن در سطح کشور با تمرکز ویژه بر تحلیل همه استان‌های کشور و شهر تهران مورد بررسی قرار گرفت. داده‌ها پس از پاک‌سازی و پیش‌پردازش، با استفاده از روش‌های تحلیل اکتشافی، بصری‌سازی و مدل جنگل تصادفی مورد تحلیل قرار گرفتند. مهم‌ترین یافته‌های این پژوهش به شرح زیر است:

\subsection{یافته‌های اصلی}

\begin{enumerate}
\item \textbf{الگوهای توزیع قیمت:} توزیع قیمت‌ها دارای چولگی مثبت است و تعداد محدودی از املاک بسیار گران‌قیمت موجب افزایش میانگین قیمت شده‌اند. همچنین پراکندگی زیاد قیمت‌ها بیانگر ناهمگنی قابل‌توجه بازار مسکن در کشور است.
                    
\item \textbf{تمرکز جغرافیایی داده‌ها:} بیش از یک‌سوم از کل داده‌ها (۷۶۶۳ از ۱۸۲۹۲ رکورد) مربوط به استان تهران است. این تمرکز بالا نشان‌دهنده حجم بالای معاملات مسکن در پایتخت و همچنین دسترسی بهتر به داده‌های این استان می‌باشد. استان‌های البرز، فارس و خوزستان در رتبه‌های بعدی قرار دارند.
    
\item \textbf{شکاف قیمتی بین استان‌ها:} بین قیمت مسکن در استان تهران و سایر استان‌ها شکاف بسیار چشمگیری وجود دارد. میانگین قیمت هر متر مربع در تهران (حدود ۱۱۱۵۰۰ تومان) بیش از دو برابر میانگین قیمت در استان دوم (فارس با حدود ۴۷۷۰۰ تومان) است. این شکاف نشان‌دهنده تمرکز شدید فعالیت‌های اقتصادی، فرصت‌های شغلی، جمعیت و تقاضای مسکن در پایتخت است.

\item \textbf{تأثیر موقعیت مکانی بر قیمت:} در سطح شهر تهران نیز تفاوت‌های قیمتی بسیار قابل‌توجهی بین مناطق مختلف وجود دارد. منطقه ۳ (شامل محله‌های ونک، جردن، قیطریه) با میانگین قیمت ۲۳۰۱۸۵ تومان گران‌ترین و منطقه ۲۰ (شهرری) با ۵۸۰۱۹ تومان ارزان‌ترین منطقه تهران است. این اختلاف حدود ۱۷۲۰۰۰ تومانی نشان‌دهنده نابرابری شدید در بازار مسکن پایتخت است.

\item \textbf{الگوهای زمانی معاملات:} تعداد معاملات مسکن در روزهای مختلف ماه الگویی نوسانی را نشان می‌دهد. این نوسانات می‌تواند با زمان وصول حقوق و تصمیم‌گیری‌های مالی خانوارها مرتبط باشد.
            
\item \textbf{همبستگی متغیرها:} به‌جز رابطه قوی بین قیمت کل و قیمت هر متر مربع، سایر متغیرهای عددی همبستگی خطی ضعیفی با هم دارند که بیانگر پیچیدگی عوامل مؤثر بر بازار مسکن است.
                
\item \textbf{عملکرد مدل پیش‌بینی:} مدل جنگل تصادفی با ضریب تعیین \lr{$R^2=0.9429$} توانست بخش عمده‌ای از تغییرات قیمت مسکن را پیش‌بینی کند که بیانگر عملکرد مناسب این مدل برای تحلیل داده‌های بازار مسکن است.
                
\item \textbf{اهمیت ویژگی‌ها:} تحلیل اهمیت ویژگی‌ها نشان داد که «شهر» با ضریب اهمیت ۵/۵۰ درصد، مهم‌ترین عامل مؤثر بر قیمت مسکن است. پس از آن، متراژ، روز معامله، سن بنا، نوع سازه و استان به‌ترتیب بیشترین نقش را در پیش‌بینی قیمت ایفا می‌کنند.
\end{enumerate}

\subsection{پیشنهادات کاربردی}

بر اساس یافته‌های این پژوهش، پیشنهادهای زیر ارائه می‌شود:

\begin{enumerate}
    \item \textbf{توسعه داده‌های استان‌های کم‌نمایش:} برای استان‌هایی که داده‌های محدودی دارند (مانند کهگیلویه و بویراحمد، یزد و کرمان)، نیاز به جمع‌آوری داده‌های بیشتر برای تحلیل دقیق‌تر وجود دارد.

    \item \textbf{تحلیل عمیق‌تر عوامل قیمت‌گذاری:} با توجه به شکاف قیمتی بالا، تحقیق در مورد عوامل اقتصادی ـ اجتماعی مؤثر بر قیمت مسکن در مناطق مختلف (مانند دسترسی به حمل‌ونقل عمومی، کاربری اراضی، کیفیت مدارس و ...) می‌تواند به درک بهتر این تفاوت‌ها کمک کند.

    \item \textbf{پایش مستمر بازار مسکن:} با توجه به نوسانات روزانه قیمت و حجم معاملات، ایجاد یک سامانه پایش مستمر بازار مسکن می‌تواند به تصمیم‌گیران در شناسایی روندها و اتخاذ سیاست‌های به‌موقع کمک کند.

    \item \textbf{تحلیل منطقه‌ای برای سرمایه‌گذاری:} سرمایه‌گذاران می‌توانند با استفاده از نتایج این پژوهش، تفاوت‌های قیمتی میان استان‌ها و مناطق مختلف را بهتر ارزیابی کرده و تصمیمات سرمایه‌گذاری خود را بر اساس شواهد آماری اتخاذ نمایند.
\end{enumerate}

\subsection{محدودیت‌ها و جهت‌گیری تحقیقات آینده}

این پژوهش دارای محدودیت‌هایی نیز می‌باشد که می‌تواند جهت‌گیری تحقیقات آینده را مشخص کند:

\begin{itemize}
        \item \textbf{نامتوازن بودن توزیع داده‌ها:} تمرکز بخش قابل‌توجهی از داده‌ها در استان تهران می‌تواند بر برخی نتایج تحلیل و مدل‌سازی اثر گذاشته باشد. افزایش حجم داده‌های سایر استان‌ها می‌تواند به تعمیم‌پذیری بهتر نتایج کمک کند.

\item \textbf{داده‌های مقطعی:} داده‌های مورد استفاده در این پژوهش مقطعی هستند و روندهای زمانی بلندمدت را پوشش نمی‌دهند. استفاده از داده‌های سری زمانی می‌تواند به شناسایی روندهای قیمتی و فصلی کمک کند.
        
    \item \textbf{عوامل کیفی:} عواملی مانند کیفیت ساخت، امکانات رفاهی ساختمان، وضعیت دسترسی و ... در این داده‌ها وجود ندارند. ادغام داده‌های کیفی می‌تواند دقت مدل‌های پیش‌بینی را افزایش دهد.
        
    \item \textbf{تأثیرات کلان اقتصادی:} بررسی تأثیر متغیرهای کلان اقتصادی مانند نرخ تورم، نرخ ارز و نرخ سود بانکی بر بازار مسکن می‌تواند به درک بهتر پویایی‌های این بازار کمک 
        کند.
\end{itemize}

\subsection{جمع‌بندی نهایی}

تحلیل داده‌های مسکن ایران نشان می‌دهد که بازار مسکن کشور با نابرابری‌های جغرافیایی عمیقی مواجه است. تمرکز شدید معاملات و قیمت‌های بالا در تهران، نیاز به سیاست‌های متوازن‌کننده دارد و لزوم توسعه متوازن مناطق مختلف کشور را بیش از پیش آشکار می‌سازد. همچنین، بر اساس نتایج مدل یادگیری ماشین، متغیر «شهر» مهم‌ترین عامل مؤثر بر قیمت مسکن است و پس از آن متراژ، روز معامله، سن بنا و نوع سازه قرار دارند. این نتایج می‌تواند در تصمیم‌گیری‌های خرید، فروش و سرمایه‌گذاری مورد استفاده قرار گیرد.

   \end{document} 